  \documentclass{article}
\usepackage{amsmath}
\usepackage{graphicx}
\graphicspath{ {images/} }
\usepackage[spanish]{babel}
\usepackage{bbm}
\usepackage{amsthm}
\usepackage{authblk}
\usepackage{hyperref}
\usepackage[utf8]{inputenc}
\usepackage{mathrsfs}
\usepackage{amsfonts}
\usepackage{amssymb}
\usepackage{enumerate}
\usepackage[left=3cm,right=2cm,top=2cm,bottom=2cm]{geometry}
\usepackage{epsfig}
 \renewcommand{\baselinestretch}{0.93}
 \thispagestyle{empty}
 \pagestyle{empty}
\setlength{\textheight}{53pc} \setlength{\textwidth} {36pc}
\setlength{\topmargin}{-4mm}
\usepackage{multicol}
\newcommand*{\email}[1]{%
    \normalsize\href{mailto:#1}{#1}\par
    }
\setlength{\columnsep}{1cm}
\newcommand{\N}{\mathbb{N}}
\newcommand{\calM}{\cal{M}}
\newcommand{\calP}{\cal{P}}
\newcommand{\F}{\mathbb{F}}
\newcommand{\R}{\mathbb{R}}
\newcommand{\Z}{\mathbb{Z}}
\newcommand{\Q}{\mathbb{Q}}
\newcommand{\C}{\mathbb{C}}
\newcommand{\bbH}{\mathbb{H}}


\theoremstyle{definition}
\newtheorem{ejer}{Ejercicio}
\author{Marcos Morales}
\affil{Universidad Finis Terrae}
\title{Primera Semana\\}



	


\begin{document}

\begin{center}
    \Large{\textbf{Taller con décimas}}
\end{center}


\setlength{\parindent}{0cm}


\textbf{Ejercicio 1:} Calcule las siguientes integrales 

\begin{align*}
    &\int_0^1\int_x^1 \frac{sin(y)}{y}\cos\left( \frac{x}{y}\right) dy dx \\
    &\int_0^1\int_0^{\sqrt{1-x^2}} e^{\sqrt{x^2+y^2}}dydx
\end{align*}





\textbf{Ejercicio 2:} Considere la región $R$ acotada por las gráficas de $y=x^3$, $y=x^2+4$ y $5x+y=0$. Realice el dibujo de la situación y después, utilizando integrales dobles determine el área de $R$. \\ \\








\textbf{Ejercicio 3:} Sea $S$ el sólido generado en el primer octante interior a la semiesfera $z=\sqrt{16-x^2-y^2}$ e interior al cilindro $x^2+y^2=4$. Usando integrales dobles calcule el volumen de $S$. \\ \\


\textbf{Ejercicio 4:} Sea $E$ la región acotada por la gráfica del plano $2x+3y+z=12$ en el primer octante. Calcule

\begin{align*}
    \iiint_E x-y+z dV
\end{align*}



\end{document} 